\chapter{Tones: Theory and Computation} 
\label{chap:tones}

Languages with tonal distinctions have been studied as a more or less separate 
subfield known as tonal phonology. This division reflects a natural difference 
between the nature of tone and that of segments, as evidenced by a large body of 
empirical evidence from Asian and African languages. Rather than being inherent 
properties of vowels, tones may change independently of the segments that bear 
them. This property challenges traditional linear rule-writing approaches, which 
provide a cumbersome means of capturing tonal generalizations. Consequently, any 
discussion of tone necessarily involves a discussion of its representations, 
ranging from accent marks and feature-based representations to phonetically 
based notational systems \citep{chao1930system}. These approaches can be found 
throughout both the African and Asian tonal phonology literature. Among them, 
Goldsmith's (1976) autosegmental theory has become the most influential and widely 
adopted framework. The first half of this chapter provides an overview of the major 
approaches to tonal representation and typology.

The issue of representation further poses important requirements for computational 
modeling. Since the introduction of autosegmental theory by 
\citet{goldsmith1976autosegmental}, there has been a clear shift in theoretical 
phonology toward structurally enriched representations. In contrast, computational 
approaches have been slower to adopt such representations. 
Works by \citet{jardine16dissertation,chandlee2019autosegmental,oakden2020computational} 
argue that linguistically meaningful representations can reduce computational complexity. 
These representations include autosegmental graphs \citep{jardine16dissertation}, 
metrical structures \tocite{lamount}, and feature-based representations for segmental 
learning \tocite{magda}. These proposals, together with their computational 
implementations, highlight an important direction for computational modeling: 
the use of linguistically informed representations to improve interpretability, 
data efficiency, and linguistic adequacy through the development of small, 
specialized language models. The second half of this chapter reviews 
computational approaches to tonal phonology, including finite-state transducers 
and other computational models. It also discusses the challenges and 
opportunities involved in developing computational models that effectively 
capture the complexities of tonal phonology.

\section{Representations of Tone}

Estimates of the proportion of the world's tonal languages vary 
from 42.9\% \tocite{thet}

\citet{yip2002tone}, for example, suggests that as many as 60--70\% of the world's 
languages may be tonal, whereas \citet{maddieson2013tone} estimates the proportion 
to be closer to 40\%. This discrepancy reflects differences in language sampling as 
well as in the criteria used to distinguish tone languages from languages with pitch 
accent or marginal tonal contrasts. Nevertheless, tone languages have a non-even 
geographical distribution: they are especially widespread in Africa, Central America, 
and East and Southeast Asia, while occurring more sparsely in Eurasia, North America, 
and Australia 
\citep{welmers1962, maddieson2013tone, kao2017phonological, 
	gordon2016phonological}.

Since tonal languages are widespread across different 
geographic regions, different conventions for representing 
tones have been developed to accommodate the languages in 
questions. \citet{yip2002tone} classified three main notational traditions 
for tones: Africanist, Asianist, and Central Americanist.
I group the latter two together for the present discussion since 
both use numerical notation, although they differ in how the numbers 
correspond to pitch levels and contours. 

It is worth mentioning that, in this section, I use the terms ``notation'' and 
``representation'' interchangeably. Although the two terms differ, all the methods 
discussed below have been used in theoretical or computational analyses of tonal 
phonology. I review five such systems: letter notation, number notation, feature-based 
representation, Q-Theory, and autosegmental representation. The first two function 
primarily as notational systems in previous work, whereas the remaining three are more 
formal representational frameworks that have also been applied beyond tonal phonology.

\subsection{Letter System}

An Africanist convention is to represent tone using accent marks. Tones are commonly 
marked with diacritics, even when they are not represented in a language’s everyday 
orthography. Table~\ref{tab:africanist} shows a basic system.
\begin{table}[ht]
	\centering
	\caption{Common Africanist tone notation}
	\label{tab:africanist}
	\begin{tabular}{lll}
		\hline
		\textbf{Tone} & \textbf{Symbol} & \textbf{Letter} \\
		\hline
		High      & \textit{á} & H      \\
		Low       & \textit{à} & L      \\
		Mid       & \textit{ā} & M      \\
		Falling   & \textit{â} & F   \\
		Rising    & \textit{ǎ} & R   \\
		%		Fall-rise & `\/\` & FR/HLH \\
		%		Rise-fall & \textit{a᷈} & RF/LHL \\
		\hline
	\end{tabular}
\end{table}

It should be noted that Africanist notation varies across linguists 
and languages. First, the unmarked tone in a two-tone system may be 
left without a diacritic; similarly, the mid tone is often unmarked 
in a three-tone system. Conversely, the macron (ā) is also used to 
indicate vowel length in some literature. 
\tocite{Provide examples of unmarked-tone conventions in several African languages} 

Second, the notation of contour tones may be ambiguous, particularly 
in languages in which contour tones have greater distributional 
flexibility. As shown in Table~\ref{tab:africanist}, a falling tone 
may be written with a circumflex on the nucleus vowel (\textit{â}), 
on a long vowel (\textit{â:}), or on one element of a complex 
nucleus (\textit{âi}, \textit{ân}). Additionally, contour tones
can emerge in the surface form as a morphological or phonological processes, even though the language does not have underlying contour form. One practice is to distinguish underlying contour from deritational contour with F vs HL and R vs LH. but this praxrixe is 
not consistnetly used. therefore, contour tones may alternatively be represented as (\textit{áà} = \textit{â:}). In 
languages that permit both concave and convex contour tones, 
circumflex and caron diacritics alone may be insufficient to 
represent more complex contours unambiguously.\tocite{Mende tonal 
notation system}

\begin{table}[ht]
	\centering
	\caption{Common Africanist tone notation}
	\label{tab:africanist_updated}
	\begin{tabular}{lll}
		\hline
		\textbf{Tone} & \textbf{Symbol} & \textbf{Letter} \\
		\hline
		High      & \textit{á} & H      \\
		Low       & \textit{à} & L      \\
		Mid       & \textit{ā} & M      \\
		Falling   & \textit{â,â:,áà} & F/HL   \\
		Rising    & \textit{ǎ} & R/LH   \\
		Fall-rise & `\/\` & FR/HLH \\
		Rise-fall & \textit{a} & RF/LHL \\
		\hline
	\end{tabular}
\end{table}

Although these rules describe the observed mappings, they obscure the generalization 
that the tonal components are preserved while the number of tone-bearing units is 
reduced. A representation that decomposes contour tones into sequences of level tones 
captures this generalization more directly.

\subsection{Phonetic-Based Notation}

In the Asianist tradition, tone is commonly represented using Chao’s five-level 
system \citep{chao1930system}. An example of this system is shown in Table~\ref{tab:cantonese-tones}, which lists the nine lexical tones of Cantonese.

\begin{table}[ht]
\centering
\caption{Lexical tones in Cantonese}
\label{tab:cantonese-tones}

\begin{tabular}{llll@{\qquad}llll}
\toprule
\multicolumn{4}{c}{Non-entering} &
\multicolumn{4}{c}{Entering} \\
\cmidrule(r){1-4}\cmidrule(l){5-8}
Tone & Value & Description & Example &
Tone & Value & Description & Example \\
\midrule
T1 & 55 & High level   & \textit{si1} `poem' &
T7 & 5  & High entering & \textit{sik7} `know' \\

T2 & 25 & High rising  & \textit{si2} `history' &
T8 & 3  & Mid entering  & \textit{sek8} `tin' \\

T3 & 33 & Mid level    & \textit{si3} `to try' &
T9 & 2  & Low entering  & \textit{sihk9} `eat' \\

T4 & 21 & Low falling  & \textit{si4} `time' &
   &    &               & \\
T5 & 23 & Low rising   & \textit{si5} `market' &
   &    &               & \\

T6 & 22 & Low level    & \textit{si6} `matter' &
   &    &               & \\
\bottomrule
\end{tabular}
\end{table}
 
Chao's system has been widely adopted in the study of Chinese languages and dialects.
The specific rules are as follows:
\begin{enumerate}[itemsep=1em, parsep=-1em]
	\item The natural pitch range of a speaker’s normal voice is divided into five levels, with 1 representing the lowest pitch and 5 the highest.
	\item A sequence of digits indicates the starting point, turning point(s), and the endpoint of a pitch contour. 
	\item A toneless or neutral-tone syllable may be represented by 0.
	\item An entering-tone syllable—that is, a syllable ending in a stop coda—is conventionally represented by a single digit.
\end{enumerate}

In Central Americanist notation, tones are represented in a similar way, but in a 
reversed order, with 1 representing the highest pitch and 5 the lowest. This system 
is used in the study of Mayan languages, such as K’iche’ and Q’eqchi’. 

This way of representing tones is effective in describing their productive
and perceptual details. However, this can also leads to its limitation
and explains why such system is more treated as a notation rather than a formal mental representation of speakers. Transcriptions of the same languages can vary greatly across
linguists due to allphonic (or allotonic) and interspeakers variability. 
Take Shanghai Chinese (also known as Shanghainese, a variaty in Wu Chinese) as an example. Shanghainese has five tones, and their tone values have had different numerical
descriptions from different sources, as shown in Table~\ref{tab:shanghai-tones}.

\begin{table}[ht]
\centering
\caption{Tone values reported across different sources}
\label{tab:shanghai-tones}
\begin{tabular}{lccccc}
\toprule
& \multicolumn{5}{c}{Tones} \\
\cmidrule(lr){2-6}
Sources & T1 & T2 & T3 & T4 & T5 \\
\midrule
JSS 1960          & 53              & 34               & 13 / 14           & 5                & 2               \\
Walton 1983       & 53              & 34               & 14                & 5                & 2               \\
Xu et al. 1981--83 & 53             & 34               & 13                & \underline{55}   & 12              \\
Xu \& Tang 1988   & 53 / \underline{52}
                  & 334 / \underline{434}
                  & 113 / \underline{223}
                  & \underline{55} / \underline{54}
                  & 12 / \underline{23} \\
T. Shen 1981 (1985)
                  & 52
                  & 334 / 34
                  & 113 / 13
                  & 4 / 5 / 7 (\underline{45})
                  & \underline{23} \\
Yuan 1983         & 42              & 34 / 434         & 24                & 5                & \underline{23}  \\
Norman 1988       & 42              & 35               & 24                & \underline{55}   & \underline{23}  \\
\bottomrule
\end{tabular}
\end{table}

Obviously, tonal processes can be generalized using this numerical system, but 
the generalization appears to involve arbitrary changes in pitch values (which is, 
as the reader progresses in Section \tocite{later}, same issues for feature system). For example, the fall-rise tone (T3) in Mandarin is consistently transcribed as 214. 
When two fall-rise tones are adjacent, the former one will undergo a sandhi process
and become a high rise. This process is formalized as 214 + 214 = 35 + 214. 
Although, this gives a clear depiction of the phonetic details, but the relationship 
between 214 and 35 appears arbitrary and does not reveal the phonological structure 
underlying the change. 

\subsection{Feature Theory}
Feature theory is another option for describing tones. H, L, R, F can be presented as [+H,-contour], [-H,-contour], [-H,+contour], [+H,+contour]. In other words, R
and F will be a natural class of [+contour] and L and R will be [-H]. However, this system is also inadequate for several reasons. 


Firstly, this theory fails to produce the correct rules for patterns that do exist, and also produces unattested patterns. Example \ref{ex:cantdo} lists two common rules when a H tone becomes a R after \{L,F\}, or a L tone turns into a F after \{H,R\}, and there is no possible way to describe the process by using [contour] feature because \{L,F\} or \{H,R\} does not form a natural class. Besides, the theory can also predict unattested patterns in Example \ref{ex:crazyrules}.

\ex \label{ex:cantdo}
H $\rightarrow$ R / \{L,F\} \_\_\\
L $\rightarrow$ F / \{H,R\} \_\_
\xe 

\ex \label{ex:crazyrules}
*[+H] $\rightarrow$ [+contour] / [-H]\_\_  \hfill equivalent to H $\rightarrow$ F / \{L,R\} \_\_\hfill(unattested)\\
*[-H] $\rightarrow$ [+contour] / [+H]\_\_\hfill   equivalent to L $\rightarrow$ R / \{H,F\} \_\_ \hfill(unattested)
\xe

As can be seen, this could be simply handled by ARs by adding an association line between L and H to the segment:

    \begin{tikzpicture}
        \node (s1) at (0, 0) {V} ; 
        \node[anchor=base] (t1) at (0,1) {H} ; 
        \node[anchor=base] (t2) at (1,1) {L} ;
        %association lines
        \draw [dotted] (s1) -- (t2) ;
        \draw (s1) -- (t1) ;
    \end{tikzpicture}
    

Another phenomenon which argues for the autosegmental representation
of tone is across-the-board effects, which refers to the effect that a tone change applies to the whole domain rather than just its local neighbor. As shown in Example \tocite{ex:shona} from Shona, a H tone prefix will trigger a series of L tone to change into H. However, if more prefixes stack together, the melody will be in tone alternation, which would be extremely difficult for feature theory to capture accurately.
\subsection{Q-Theory}

\citet{inkelas2013abc+} introduced \textit{Q-theory}, initially designed to account
for complex diphthongs and later expanded to contour tones. Under this
framework, each tone (or segment) can be divided into three temporally ordered,
quantized subsegments. A phonological representation follows the format
Q($q_1$, $q_2$, $q_3$), where a traditional representation, or a segment, can be
decomposed into three subsegments $q$. Here, ``q'' is loosely based on the concept
of quantized temporal subphases of the unit phonologists call a ``segment.''
A vowel with a rising-falling tone can be represented as
V(à\textsuperscript{1}á\textsuperscript{2}à\textsuperscript{3}).
\citet{inkelas2013abc+} argue that Q-theory can compensate for the difficulty of
Aperture Theory, which limits the number of aperture positions to a maximum of
two.

Q-theory is motivated by the typological observation that most complex tones do
not exceed three tonal levels and that four-tone contours are very rare
\citep{shih2013subsegmental}. It provides a useful harmony--disharmony typology
based on whether the target of change is the whole segment or a subsegment.
% See Table~\ref{tab:tone-harmony}.


% \begin{table}[htbp]
% 	\centering
% 	\caption{Tone harmony and disharmony involving whole segments and subsegments}
% 	\label{tab:tone-harmony}
% 	\begin{tabular}{lll}
% 		\hline
% 		& \textbf{Whole segment: Q} & \textbf{Subsegment: q} \\
% 		\hline
		
% 		\textbf{Harmony}
% 		&
% 		\begin{tabular}[t]{@{}l@{}}
% 			\textbf{a. Changzhi} \\
% 			(Duanmu 1994; a.o.) \\
% 			/kuə213-təʔ535/ $\rightarrow$ [kuə213-təʔ213] \\
% 			`pan, dim.'
% 		\end{tabular}
% 		&
% 		\begin{tabular}[t]{@{}l@{}}
% 			\textbf{c. Yoruba} \\
% 			(Akinbiyi \& Liberman 2001) \\
% 			/rárà/ $\rightarrow$ [rárâ] \\
% 			`elegy'
% 		\end{tabular}
% 		\\[2ex]
		
% 		\textbf{Disharmony}
% 		&
% 		\begin{tabular}[t]{@{}l@{}}
% 			\textbf{b. Pingyao} \\
% 			(Chen 1992; Yip 2002) \\
% 			/hai35 bing35/ $\rightarrow$ [hai53 bing35] \\
% 			`become ill'
% 		\end{tabular}
% 		&
% 		\begin{tabular}[t]{@{}l@{}}
% 			\textbf{d. Tianjin} \\
% 			(Chen 1985; a.o.) \\
% 			/feiL-jiL/ $\rightarrow$ [feiLH.jiL] \\
% 			`airplane'
% 		\end{tabular}
% 		\\
		
% 		\hline
% 	\end{tabular}
% \end{table}

The first limitation of Q-theory is that, even though the typological observation
suggests a lack of four-tone contours, there is in fact a typological category
found in some Chinese dialects known as the ``double circumflex''
\citep{zhu2012dipping}. Double-circumflex tones have two perceptually salient turning
points. Three variants of the double-circumflex tone can be distinguished by the
relative heights of these turning points, as shown in
Table~\ref{tab:double-circumflex}.

\begin{table}[htbp]
	\centering
	\caption{Three variants of double-circumflex tones}
	\label{tab:double-circumflex}
	\begin{tabular}{llll}
		\hline
		\textbf{Type} & \textbf{Contour} & \textbf{Variety} & \textbf{Example} \\
		\hline
		Back-low     & 4232 & Yuejiashan Mandarin (Zhongyang)
		& \textit{xiang} `make a sound' \\
		Equal-height & 3232 & Deqing Wu (Zhejiang)
		& \textit{ma} `hemp' \\
		Front-low    & 3242 & Qiyang Xiang (Hunan)
		& \textit{dong} `freeze' \\
		\hline
	\end{tabular}
\end{table}

One possible solution for accommodating these tones within Q-theory is to
consider a single contour (a fall or rise) as the subsegmental unit.
However, if a language with double-circumflex tones also has level tones (such
as 55) or short tones (also known as ``entered tone'', typically noted as 5), Q-theory is again insufficient to account
for all these tone types within a unified representation.

Secondly, Q-theory could also be seen as a flat-string version of a nonlinear
representation. Q could be seen as a tonal node, while its subsegments could be
seen as tone units. Whether a subsegment or the whole segment undergoes a change
also corresponds to the two tonal models proposed by \citet{yip1980}
 and
\citet{bao1999structure}. \citet{jardine2021qtheory} argues that Q-theory is
logically equivalent to ARs under model theory and logical interpretations. This
means that, using a set of logical formulas, there is a bijective transduction
from one representation to the other, and vice versa. The differences are
presentational, but the representations are notationally equivalent.
\citet{oakdenNotationalEquivalenceTonal2020} provides a more detailed discussion
of notational equivalence.






\subsection{Autosegmental Theory}


